% =====================================================================
% config/style.tex
% =====================================================================

% =====================================================================
% 1. GLOBAL FONT & TYPOGRAPHY
% =====================================================================
\renewcommand{\familydefault}{\rmdefault}


% =====================================================================
% 2. TABLE OF CONTENTS (TOC) FORMATTING & SPACING
% =====================================================================
\setcounter{tocdepth}{2}
\setcounter{secnumdepth}{3}

\renewcommand{\cfttoctitlefont}{\rmfamily\Large\bfseries}

% Enable dotted leader lines for main sections
\renewcommand{\cftsecleader}{\cftdotfill{\cftdotsep}}

% Set indentations and wide number box dimensions
\cftsetindents{section}{0em}{2.5em}
\cftsetindents{subsection}{2.0em}{2.5em}
\cftsetindents{subsubsection}{4.0em}{4.5em}

% Vertical spacing between items
\setlength{\cftbeforesecskip}{0.5em}
\setlength{\cftbeforesubsecskip}{0.3em}
\setlength{\cftbeforesubsubsecskip}{0.3em}

% Reset box overrides to prevent overlapping
\renewcommand{\cftsecpresnum}{}
\renewcommand{\cftsecaftersnum}{}
\renewcommand{\cftsubsecpresnum}{}
\renewcommand{\cftsubsecaftersnum}{}

% Font styling
\renewcommand{\cftsecfont}{\rmfamily\bfseries}
\renewcommand{\cftsecpagefont}{\rmfamily\bfseries}
\renewcommand{\cftsubsecfont}{\rmfamily}
\renewcommand{\cftsubsecpagefont}{\rmfamily}
\renewcommand{\cftsubsubsecfont}{\rmfamily}
\renewcommand{\cftsubsubsecpagefont}{\rmfamily}

% =====================================================================
% 3. RUNNING HEADERS AND FOOTERS
% =====================================================================

\setlength{\headheight}{14pt}

\pagestyle{fancy}
\fancyhf{}

% Header (\headertext and \currentproject are set in commands.tex;
% \currentproject changes at the start of each project in the book)
\fancyhead[L]{\rmfamily\footnotesize\scshape \headertext}
\fancyhead[R]{}

% Footer
\fancyfoot[L]{\rmfamily\footnotesize \currentproject}
\fancyfoot[R]{\rmfamily\footnotesize Page\,|\,\thepage}

% Header and footer rules
\renewcommand{\headrulewidth}{0.4pt}
\renewcommand{\footrulewidth}{0.4pt}

% Plain page style (used on section/chapter starts)
\fancypagestyle{plain}{%
  \fancyhf{}
  \fancyhead[L]{\rmfamily\footnotesize\scshape \headertext}
  \fancyhead[R]{}
  \fancyfoot[L]{\rmfamily\footnotesize \currentproject}
  \fancyfoot[R]{\rmfamily\footnotesize Page\,|\,\thepage}
  \renewcommand{\headrulewidth}{0.4pt}
  \renewcommand{\footrulewidth}{0.4pt}
}


% =====================================================================
% 4. BODY TEXT
% =====================================================================

\setlength{\parindent}{0pt}
% ragged2e's \justifying (below) resets \parindent to \JustifyingParindent,
% which defaults to 15pt; set it to 0 too so no paragraph is indented.
\setlength{\JustifyingParindent}{0pt}
\setlength{\parskip}{6pt plus 1pt minus 1pt}
\linespread{1.05}
\raggedbottom
\justifying
\sloppy
\emergencystretch=2em


% =====================================================================
% 5. COLOUR DEFINITIONS
% =====================================================================

\definecolor{HeaderGray}{RGB}{89,89,89}
\definecolor{RowGray}{RGB}{242,242,242}
\definecolor{RuleGray}{RGB}{128,128,128}
\definecolor{CodeBg}{RGB}{248,248,248}
\definecolor{CodeFrame}{RGB}{170,170,170}
\definecolor{CodeKeyword}{RGB}{0,0,153}
\definecolor{CodeComment}{RGB}{95,140,95}
\definecolor{CodeString}{RGB}{163,21,21}
\definecolor{ProfileBlue}{RGB}{31,58,110}
\definecolor{CoverNavy}{RGB}{18,78,111}
\definecolor{CoverBlue}{RGB}{20,82,119}
\definecolor{CoverGray}{RGB}{75,75,75}
\definecolor{CoverAccent}{HTML}{5EA0F5}


% =====================================================================
% 6. TABLE STYLING
% =====================================================================

\renewcommand{\arraystretch}{1.25}
\setlength{\tabcolsep}{6pt}
\arrayrulecolor{RuleGray}

% p{} aligns cell content to the top; m{} (array pkg) centers it
% vertically instead, so a short one-line cell (e.g. "Python") lines up
% with the middle of a taller multi-line cell beside it, not its top.
\renewcommand\tabularxcolumn[1]{m{#1}}
\newcolumntype{L}[1]{>{\raggedright\arraybackslash}m{#1}}
\newcolumntype{C}[1]{>{\centering\arraybackslash}m{#1}}
\newcolumntype{R}[1]{>{\raggedleft\arraybackslash}m{#1}}
\newcolumntype{Y}{>{\raggedright\arraybackslash}X}
\newcolumntype{Z}{>{\centering\arraybackslash}X}


% =====================================================================
% 7. FIGURE & TABLE FLOAT DEFAULTS
% =====================================================================

% Images use explicit per-project paths (projects/<slug>/...), so two
% projects can both contain a file called "system-architecture.png".

\renewcommand{\topfraction}{0.85}
\renewcommand{\bottomfraction}{0.65}
\renewcommand{\textfraction}{0.10}
\renewcommand{\floatpagefraction}{0.70}
\setcounter{topnumber}{3}
\setcounter{bottomnumber}{2}
\setcounter{totalnumber}{4}


% =====================================================================
% 8. SECTION HEADING FORMATTING
% =====================================================================

\titleformat{\section}
  {\rmfamily\Large\bfseries}
  {\thesection}
  {1em}
  {}

\titleformat{\subsection}
  {\rmfamily\large\bfseries}
  {\thesubsection}
  {1em}
  {}

\titleformat{\subsubsection}
  {\rmfamily\normalsize\bfseries}
  {\thesubsubsection}
  {1em}
  {}

% Tighter gaps around headings: space-before, space-after-heading,
% space-before-following-paragraph (last argument is unused here since
% titles are not run-in).
\titlespacing*{\section}{0pt}{12pt plus 2pt minus 2pt}{4pt}
\titlespacing*{\subsection}{0pt}{10pt plus 2pt minus 2pt}{3pt}
\titlespacing*{\subsubsection}{0pt}{8pt plus 1pt minus 1pt}{2pt}

% =====================================================================
% 9. FLOAT SPACING
% =====================================================================
% Tighter gap between body text and an in-text ([H]) figure than the
% article-class default, which otherwise stacks with \parskip and reads
% as an oversized gap before every diagram/screenshot.
\setlength{\intextsep}{4pt plus 1pt minus 1pt}
\setlength{\textfloatsep}{6pt plus 1pt minus 1pt}
\setlength{\floatsep}{4pt plus 1pt minus 1pt}